\section*{الفصل \ref{ch:b2:multint} --- التكاملات المنحنية والتكاملات المتعددة}

\begin{solution}{exo:b2:multint:1}
(1) القطعة $\gamma(t) = (t, t)$، $t \in [0,1]$:
\[
\int_\gamma y^2\dd x + x\dd y
= \int_0^1 (t^2 + t)\,\dd t = \frac13 + \frac12 = \frac56 .
\]
(2) القطع المكافئ $\gamma(t) = (t, t^2)$:
\[
\int_0^1 \bigl(t^4\cdot 1 + t\cdot 2t\bigr)\dd t
= \frac15 + \frac23 = \frac{13}{15} .
\]
والقيمتان مختلفتان، ومنه فإن التكامل يتعلق بالمسار: أي إن الصورة
\emph{ليست} تامة --- وبانسجام مع ذلك، $P_y = 2y \neq 1 = Q_x$، ومنه فهي ليست مغلقة
أصلا.
\end{solution}

\begin{solution}{exo:b2:multint:2}
$P = 2xy + y^3$، $Q = x^2 + 3xy^2 + 1$: $P_y = 2x + 3y^2 = Q_x$، فهي مغلقة على $\R^2$. ونبحث عن $f$ يحقق
$f_x = P$: $f = x^2y + xy^3
+ g(y)$؛ ثم يفرض $f_y = x^2 + 3xy^2 + g'(y) = Q$ أن $g'(y) = 1$، ولنأخذ $g(y) = y$. ومنه
\[
f(x, y) = x^2y + xy^3 + y
\]
\omterm{def:b2:multint:exact}{كمون} (والمجموعة $\R^2$ نجمية الشكل، ومنه كان لا بد أن يوجد \omterm{def:b2:multint:exact}{كمون}
بمبرهنة بوانكاريه --- لكن إظهاره أسرع). وبحسب \cref{thm:b2:multint:ftc}، من أجل أي قوس
من $(0,0)$ إلى $(1,2)$:
\[
\int_\gamma\omega = f(1, 2) - f(0, 0) = 2 + 8 + 2 = 12 .
\]
\end{solution}

\begin{solution}{exo:b2:multint:3}
المنطقة هي $\{0 \leq x \leq 1,\ x^2 \leq y \leq x\}$.
وبمكاملة $y$ أولا:
\[
\int_0^1\!\int_{x^2}^{x}(x + y)\,\dd y\,\dd x
= \int_0^1\Bigl(x(x - x^2)
+ \frac{x^2 - x^4}{2}\Bigr)\dd x
= \int_0^1\Bigl(\frac{3x^2}{2} - x^3 - \frac{x^4}{2}\Bigr)\dd x
= \frac12 - \frac14 - \frac1{10} = \frac{3}{20} .
\]
وبمكاملة $x$ أولا: المقطع عند الارتفاع $y \in [0, 1]$ هو $y \leq x \leq
\sqrt y$، ومنه
\[
\int_0^1\!\int_{y}^{\sqrt y}(x + y)\,\dd x\,\dd y
= \int_0^1\Bigl(\frac{y - y^2}{2}
+ y(\sqrt y - y)\Bigr)\dd y
= \frac14 - \frac16 + \frac25 - \frac13 = \frac{3}{20} .
\]
\end{solution}

\begin{solution}{exo:b2:multint:4}
\emph{التكامل الأول.} على القرص $D_R$، بالإحداثيات القطبية:
\[
\iint_{D_R}\frac{\dd x\,\dd y}{(1 + x^2 + y^2)^2}
= \int_0^{2\pi}\!\!\int_0^R
\frac{\rho\,\dd\rho\,\dd\alpha}{(1 + \rho^2)^2}
= 2\pi\Bigl[-\frac{1}{2(1 + \rho^2)}\Bigr]_0^R
= \pi\Bigl(1 - \frac{1}{1 + R^2}\Bigr)
\xrightarrow[R \to \infty]{} \pi .
\]
\emph{التكامل الثاني.} ربع القرص هو $0 \leq \alpha \leq
\frac\pi2$، $0 \leq \rho \leq 1$، و$xy =
\rho^2\cos\alpha\sin\alpha$:
\[
\iint_{D'}xy\,\dd x\,\dd y
= \int_0^{\pi/2}\!\!\cos\alpha\sin\alpha\,\dd\alpha
\int_0^1 \rho^3\,\dd\rho
= \frac12\cdot\frac14 = \frac18 .
\]
\end{solution}

\begin{solution}{exo:b2:multint:5}
بحسب \cref{cor:b2:multint:area} مع $x = \cos^3 t$، $y = \sin^3 t$: $x' = -3\cos^2 t\sin t$، $y' = 3\sin^2 t\cos t$، ومنه
\[
xy' - yx' = 3\cos^4 t\sin^2 t + 3\sin^4 t\cos^2 t
= 3\sin^2 t\cos^2 t = \frac{3}{4}\sin^2 2t
= \frac{3}{8}(1 - \cos 4t) .
\]
ومنه
\[
A = \frac12\int_0^{2\pi}\frac38(1 - \cos 4t)\,\dd t
= \frac{3}{16}\cdot 2\pi = \frac{3\pi}{8} .
\]
(\omterm{pb:b2:curves:1}{والنجمية} محتواة في قرص الوحدة الذي مساحته $\pi$؛ وثلاثة أثمان
$\pi$ قيمة معقولة لشكلها النجمي ذي نقط الرجوع الأربع.)
\end{solution}

\begin{solution}{exo:b2:multint:6}
\emph{بالتكديس:} فوق كل $(x, y)$ من قرص الوحدة $D$، يجري $z$ من
$x^2 + y^2$ إلى $1$:
\[
V = \iint_D \bigl(1 - x^2 - y^2\bigr)\dd x\,\dd y
= \int_0^{2\pi}\!\!\int_0^1 (1 - \rho^2)\rho\,\dd\rho\,\dd\alpha
= 2\pi\Bigl(\frac12 - \frac14\Bigr) = \frac\pi2 .
\]
\emph{بالتشريح:} المقطع عند الارتفاع $z \in [0, 1]$ هو القرص $x^2 + y^2 \leq z$، ومساحته
$\pi z$:
\[
V = \int_0^1 \pi z\,\dd z = \frac\pi2 .
\]
\end{solution}

\begin{solution}{exo:b2:multint:7}
بالإحداثيات الكروية يكون عنصر الحجم
$r^2\cos\varphi\,\dd r\,\dd\theta\,\dd\varphi$
(\cref{ex:b2:multint:spherical})، ويتكامل الجزء الزاوي إلى $4\pi$ ($2\pi$ من $\theta$، $\int_{-\pi/2}^{\pi/2}\cos = 2$).
وحجم القشرة هو
\[
\int_a^b 4\pi r^2\,\dd r = \frac43\pi\bigl(b^3 - a^3\bigr).
\]
أما التكامل الثاني، فالدالة تحت التكامل $1/r$ لا تتعلق إلا بالمقدار
$r$:
\[
\iiint_B \frac{\dd x\,\dd y\,\dd z}{r}
= \int_0^R 4\pi r^2\cdot\frac1r\,\dd r
= 4\pi\,\frac{R^2}{2} = 2\pi R^2 .
\]
(وتنفجر الدالة تحت التكامل في المبدأ، لكن دون ضرر: فالمقدار $r^2/r = r$
\omterm{def:b2:metric:continuity}{متصل} --- ويتقارب التكامل على القشور $\varepsilon \leq r
\leq R$ عندما $\varepsilon \to 0$، وهذا هو
المعنى الدقيق للنص. وهذا النوع من الحساب هو الخطوة الأولى نحو مبرهنة
نيوتن القائلة إن كرة مصمتة متجانسة تجذب كنقطة كتلة في مركزها.)
\end{solution}

\begin{solution}{exo:b2:multint:8}
الدالة تحت التكامل $g(t; x, y) = xP(tx, ty) + yQ(tx, ty)$ من الصنف $\mathcal{C}^1$ بدلالة $(x, y)$، ومتصلة
بدلالة $t$، ومشتقاتها الجزئية متصلة على $[0,1] \times U$؛ وتعطي المفاضلة
تحت علامة التكامل (\cref{ch:b2:integration}، مطبقة على الفترة المتراصة $[0,1]$ للمتغير
$t$، فالهيمنة تلقائية هناك)
\[
f_x(x, y)
= \int_0^1 \bigl(P(tx, ty) + tx\,P_x(tx, ty)
+ ty\,Q_x(tx, ty)\bigr)\dd t .
\]
وباستعمال الانغلاق $Q_x = P_y$:
\[
tx\,P_x(tx, ty) + ty\,P_y(tx, ty)
= t\,\frac{\dd}{\dd t}\bigl[P(tx, ty)\bigr] ,
\]
ومنه فإن الدالة تحت التكامل هي $P(tx, ty) + t\frac{\dd}{\dd t}P(tx, ty) =
\frac{\dd}{\dd t}\bigl[t\,P(tx, ty)\bigr]$ و
\[
f_x(x, y) = \Bigl[t\,P(tx, ty)\Bigr]_0^1 = P(x, y) .
\]
وبالتناظر $f_y = Q$ (بالحساب نفسه مع استعمال $P_y = Q_x$ في الاتجاه الآخر).
ولاحظ أين تدخل الفرضية: فالمقدار $f$ معرَّف بالمكاملة على طول
القطعة $[0, M]$، وهي تقع في $U$ بالضبط لأن $U$ نجمي الشكل.
\end{solution}

\begin{solution}{exo:b2:multint:9}
على الشريط $[0, A] \times [0, \infty)$ لا تكون الدالة $(x, y)
\mapsto e^{-xy}\sin x$ قابلة للمكاملة مطلقا حتى
$y =
\infty$ بانتظام بالمعنى الساذج، لكن كل تكامل مكرر يتقارب وينتج تساويهما
من مبرهنة فوبيني على $[0, A]
\times [0, B]$ مضافا إليها نهاية $B \to \infty$ (فالذيل $\int_0^A\int_B^\infty e^{-xy}\abs{\sin x}\,\dd y\,\dd x \leq
\int_0^A \frac{e^{-Bx}\abs{\sin x}}{x}\dd x \leq \int_0^A e^{-Bx}
\dd x\to 0$،
باستعمال $\abs{\sin x} \leq x$).

\emph{بمكاملة $y$ أولا:} $\int_0^\infty e^{-xy}\,\dd y = \frac1x$ من أجل $x
> 0$، ومنه فإن التكامل
الأول هو $\int_0^A \frac{\sin x}{x}\,\dd x$.

\emph{بمكاملة $x$ أولا:} تعطي مكاملتان بالتجزئة (أو أخذ الجزء
التخيلي للمقدار $\int_0^A e^{(i - y)x}\dd x$)
\[
\int_0^A e^{-xy}\sin x\,\dd x
= \frac{1 - e^{-Ay}(y\sin A + \cos A)}{1 + y^2} .
\]
وبالمكاملة بدلالة $y$ على $[0, \infty)$، ينفصل الحد $\int_0^\infty
\frac{\dd y}{1 + y^2} = \frac\pi2$:
\[
\int_0^A \frac{\sin x}{x}\,\dd x
= \frac{\pi}{2}
- \int_0^\infty e^{-Ay}\,\frac{y\sin A + \cos A}{1 + y^2}\,\dd y .
\]
والباقي محصور بالمقدار $\int_0^\infty e^{-Ay}\frac{y +
1}{1 + y^2}\dd y \leq \int_0^\infty e^{-Ay}\cdot\frac{1+y}{1+y^2}
\,\dd y \to 0$ عندما $A \to \infty$ (بالتقارب المهيمن، أو
بالحصر الفج $\frac{1 + y}{1 + y^2} \leq \frac32$ الذي يعطي $\frac{3}{2A}$). ومنه $\int_0^\infty\frac{\sin x}{x}\dd x =
\frac\pi2$ --- وهي القيمة
نفسها المحصل عليها في \cref{ch:b2:integration} بمفاضلة تكامل ذي وسيط؛ وهنا تقوم
مبرهنة فوبيني بالعمل بدلا من ذلك.
\end{solution}

\begin{solution}{exo:b2:multint:10}
نوسّم \omterm{def:b2:curves:length}{بطول القوس} $s \in [0, 2\pi]$، ومنه $x'^2 + y'^2 = 1$، وننسحب بحيث يكون $\int_0^{2\pi} x(s)\,\dd s = 0$.
وبحسب \cref{cor:b2:multint:area}،
\[
2A = \oint x\,\dd y - y\,\dd x
= \int_0^{2\pi}\bigl(xy' - yx'\bigr)\dd s .
\]
وبمكاملة $\oint y\,\dd x$ بالتجزئة على الدور (فالحدود الحدية تتلاشى بالدورية)،
$-\int yx' = \int y'x$، ومنه فإن $2A = 2\int_0^{2\pi}xy'\,\dd s$ فعلا. ثم يعطي $2xy' \leq x^2 + y'^2$
\[
2A \leq \int_0^{2\pi}\bigl(x^2 + y'^2\bigr)\dd s
= \int_0^{2\pi} x^2 + \int_0^{2\pi}\bigl(1 - x'^2\bigr)
= 2\pi - \int_0^{2\pi}\bigl(x'^2 - x^2\bigr)\dd s .
\]
وتجعل متراجحة فيرتنغر (من تمارين \cref{ch:b2:fourier}: من أجل دالة دورية بالدور
$2\pi$ من الصنف $\mathcal{C}^1$ ومتوسطها معدوم، $\int x^2 \leq \int x'^2$) التكامل الأخير
موجبا: $A \leq \pi$. ويستلزم التساوي التساوي في متراجحة فيرتنغر ($x(s) =
a\cos s + b\sin s$)
وفي $2xy' \leq x^2 + y'^2$ ($y' = x$ نقطيا)، وهو ما يفرض $y = a\sin s - b\cos s + c$: فالمنحنى دائرة
الوحدة (بتمركز مناسب). وبما أن منحنى طوله $L$ يُحاكى إلى الطول
$2\pi$، يكون النص العام $A \leq
\frac{L^2}{4\pi}$: أي إن الدائرة، من بين جميع
المنحنيات المغلقة ذات المحيط المعطى، تحصر أكبر \omterm{def:b2:multint:domain}{مساحة}.
\end{solution}

\begin{solution}{exo:b2:multint:11}
بالإحداثيات الكروية $z = r\sin\varphi$ و$\dd x\,\dd
y\,\dd z = r^2\cos\varphi\,\dd r\,\dd\theta\,\dd\varphi$:
\[
\iiint_B z^2
= \int_0^R r^4\,\dd r\int_0^{2\pi}\dd\theta
\int_{-\pi/2}^{\pi/2}\sin^2\varphi\cos\varphi\,\dd\varphi
= \frac{R^5}5\cdot2\pi\cdot
\Bigl[\frac{\sin^3\varphi}3\Bigr]_{-\pi/2}^{\pi/2}
= \frac{4\pi R^5}{15}.
\]
وبتناظر الكرة المصمتة عند تبديل الإحداثيات، $\iiint_B x^2 = \iiint_B y^2 = \iiint_B z^2$، ومنه $\iiint_B(x^2
+ y^2 + z^2) = 3\cdot\frac{4\pi R^5}{15} = \frac{4\pi R^5}5$.
وللتحقق بالقشور: $\int_0^R r^2\cdot 4\pi r^2\,\dd r = \frac{4\pi
R^5}5$ --- فالدالة تحت التكامل $r^2$ ثابتة على
الكرة ذات نصف القطر $r$، ومساحتها $4\pi r^2$.
\end{solution}

\begin{solution}{exo:b2:multint:12}
نوسّم الضلع من $(x_i, y_i)$ إلى $(x_{i+1}, y_{i+1})$ بالعلاقة $\gamma(t) = \bigl((1-t)x_i + tx_{i+1},\ (1-t)y_i +
ty_{i+1}\bigr)$. ومساهمته في
$\frac12\oint(x\,\dd y - y\,\dd x)$ هي
\[
\frac12\int_0^1\Bigl(\bigl((1-t)x_i +
tx_{i+1}\bigr)(y_{i+1} - y_i) - \bigl((1-t)y_i +
ty_{i+1}\bigr)(x_{i+1} - x_i)\Bigr)\dd t ,
\]
وبما أن $\int_0^1\bigl((1-t)u + tv\bigr)\dd t = \frac{u +
v}2$، فإن هذا يساوي
\[
\frac14\Bigl((x_i + x_{i+1})(y_{i+1} - y_i) - (y_i +
y_{i+1})(x_{i+1} - x_i)\Bigr)
= \frac12\bigl(x_iy_{i+1} - x_{i+1}y_i\bigr),
\]
مع تلاشي الحدود المتقاطعة. وبالجمع على الأضلاع وعددها $m$ نحصل
على صيغة رباط الحذاء، بحسب \cref{cor:b2:multint:area}. والمثلث $(0,0), (1,0), (0,1)$: $\frac12\bigl((0\cdot0 - 1\cdot0) +
(1\cdot1 - 0\cdot0) + (0\cdot0 - 0\cdot1)\bigr) = \frac12$، وهي
\omterm{def:b2:multint:domain}{المساحة} الصحيحة.
\end{solution}

\begin{solution}{pb:b2:multint:1}
\textbf{1.} تُوصف الكرة المصمتة $B_n(R)$ بالحصور المكررة $-R \leq x_n \leq R$، ثم
$\abs{x_{n-1}} \leq \sqrt{R^2 -
x_n^2}$، وهكذا؛ والتعويض $x_i = Ru_i$ في كل من التكاملات $n$ ذات
المتغير الواحد يضرب كلا منها في $R$ ويرسل الحدود على حدود $B_n(1)$:
$v_n(R) = R^n\,v_n(1) =
V_nR^n$.

\textbf{2.} بالتشريح على طول $x_n = t$: مقطع $B_n(1)$ هو الكرة المصمتة
$B_{n-1}\bigl(\sqrt{1 - t^2}\bigr)$، ومنه، بحسب السؤال 1،
\[
V_n = \int_{-1}^1 v_{n-1}\bigl(\sqrt{1 - t^2}\bigr)\,\dd t
= V_{n-1}\int_{-1}^{1}(1 - t^2)^{\frac{n-1}2}\,\dd t .
\]

\textbf{3.} مع $t = \sin\theta$، $\dd t =
\cos\theta\,\dd\theta$ و$(1 - t^2)^{\frac{n-1}2} =
\cos^{n-1}\theta$ على $\intcc{-\pi/2}{\pi/2}$:
\[
\int_{-1}^1(1 - t^2)^{\frac{n-1}2}\dd t
= \int_{-\pi/2}^{\pi/2}\cos^n\theta\,\dd\theta
= 2\int_0^{\pi/2}\cos^n\theta\,\dd\theta = 2W_n,
\]
والمقدار $\theta \mapsto \frac\pi2 - \theta$ يبادل بين صيغتي الجيب وجيب التمام للمقدار $W_n$.

\textbf{4.} نكتب $\sin^n = \sin^{n-2}(1 - \cos^2)$ ونكامل $\int\sin^{n-2}\cos\cdot\cos$ بالتجزئة ($v =
\frac{\sin^{n-1}}{n-1}$):
\[
W_n = W_{n-2} - \frac{W_n}{n-1}
\quad\Longrightarrow\quad
W_n = \frac{n-1}{n}W_{n-2}.
\]
ومنه $nW_nW_{n-1} = (n-1)W_{n-1}W_{n-2}$: أي إن المتتالية $(nW_nW_{n-1})$ ثابتة وتساوي $1\cdot W_1W_0 =
1\cdot\frac\pi2$، ومنه
$W_nW_{n-1} = \frac{\pi}{2n}$.

\textbf{5.} يعطي السؤالان 2 و3 المقدار $V_n = 2W_nV_{n-1}$، مرتين:
\[
V_n = 2W_n\cdot 2W_{n-1}\,V_{n-2}
= 4\,\frac{\pi}{2n}\,V_{n-2} = \frac{2\pi}n\,V_{n-2}.
\]

\textbf{6.} انطلاقا من $V_0 = 1$: $V_{2k} = \frac{2\pi}{2k}V_{2k-2}
= \frac\pi kV_{2k-2}$، ومنه $V_{2k} = \frac{\pi^k}{k!}$ بالتراجع.
وانطلاقا من $V_1 = 2$: $V_{2k+1} =
\frac{2\pi}{2k+1}V_{2k-1}$، ومنه
\[
V_{2k+1} = 2\prod_{j=1}^k\frac{2\pi}{2j+1}
= \frac{2^{k+1}\pi^k}{1\cdot3\cdot5\cdots(2k+1)} .
\]

\textbf{7.} $V_1 = 2$، $V_2 = \pi \approx 3.142$، $V_3 =
\frac{4\pi}3 \approx 4.189$، $V_4 = \frac{\pi^2}2 \approx
4.935$، $V_5 = \frac{8\pi^2}{15} \approx 5.264$، $V_6 =
\frac{\pi^3}6 \approx 5.168$،
$V_7 = \frac{16\pi^3}{105}
\approx 4.725$. والنسبة ذات الخطوة الواحدة هي $V_n/V_{n-1} = 2W_n$، و$(W_n)$ متناقصة
($\sin^n \leq \sin^{n-1}$ نقطيا). والآن $2W_5 = 2\cdot\frac45\cdot\frac23 =
\frac{16}{15} > 1$ بينما $2W_6 =
2\cdot\frac56\cdot\frac34\cdot\frac12\cdot\frac\pi2 =
\frac{5\pi}{16} < 1$: فالنسب تتجاوز $1$
حتى $n = 5$ وتقل عن $1$ ابتداء من $n = 6$ --- ومنه فإن $(V_n)$
متزايدة حتى قيمتها العظمى $V_5$ ثم متناقصة.

\textbf{8.} من أجل $n \geq 13 > 4\pi$: $V_n/V_{n-2} = 2\pi/n <
\tfrac12$، ومنه $V_{n} \leq C\cdot 2^{-n/2}$ بثابت ثابت؛ وأفضل
من ذلك، من أجل أي $q > 0$، يكون $2\pi/n < q^2$ من أجل $n$ كبيرة، ومنه
$V_n/q^n \to 0$: فالتلاشي يتفوق على كل متتالية هندسية. وينتج تقارب $\sum V_n$
من النسبة $V_n/V_{n-2} \to 0$ (بالمقارنة مع متسلسلة هندسية ابتداء من رتبة ما).

\textbf{9.} $\sum_{k\geq0}V_{2k}x^{2k} =
\sum_{k\geq0}\frac{(\pi x^2)^k}{k!} = \eu^{\pi x^2}$، وهي المتسلسلة الأسية (\cref{ch:b2:powerseries})، وهي متقاربة من
أجل كل $x$. وعند $x = 1$: $\sum_kV_{2k} = \eu^\pi \approx
23.14$.

\textbf{10.} على المكعب $\intcc{-R}R^n$ تكون الدالة تحت التكامل هي الجداء
$\prod_i\eu^{-x_i^2}$، ومنه يتفكك التكامل المكرر: $\bigl(\int_{-R}^R\eu^{-t^2}\dd t\bigr)^n$. وبجعل $R \to \infty$ واستعمال
$\int_\R\eu^{-t^2}\dd t = \sqrt\pi$
(\cref{ex:b2:multint:polar}): $I_n = \pi^{n/2}$.

\textbf{11.} يعطي $t = u^2$ المقدار $\Gamma(\tfrac12) =
\int_0^\infty t^{-1/2}\eu^{-t}\dd t =
2\int_0^\infty\eu^{-u^2}\dd u = \sqrt\pi$. وبتكرار $\Gamma(s+1) = s\Gamma(s)$: $\Gamma(k+1) = k!\,\Gamma(1) =
k!$، و
\[
\Gamma\Bigl(k + \frac32\Bigr)
= \Bigl(k + \frac12\Bigr)\Bigl(k - \frac12\Bigr)\cdots
\frac12\cdot\Gamma\Bigl(\frac12\Bigr)
= \frac{(2k+1)(2k-1)\cdots1}{2^{k+1}}\,\sqrt\pi .
\]

\textbf{12.} نضع $F_n = \pi^{n/2}/\Gamma(\frac n2 + 1)$. وبما أن $\Gamma(\frac n2 + 1) = \frac n2\,\Gamma(\frac n2) =
\frac n2\,\Gamma(\frac{n-2}2 + 1)$، نحصل على $F_n =
\frac{2\pi}nF_{n-2}$: أي
التراجع نفسه للمقدار $V_n$ (السؤال 5). والحالتان الأساسيتان:
$F_1 = \sqrt\pi/\Gamma(\frac32) =
\sqrt\pi/(\frac{\sqrt\pi}2) = 2 = V_1$ و$F_2 =
\pi/\Gamma(2) = \pi = V_2$. وبالتراجع $V_n = F_n$ من أجل كل $n$؛ ويعيد السؤال
11 تحويل ذلك إلى الصيغتين المغلقتين في السؤال 6.

\textbf{13.} مع $r = \sqrt t$: $\int_0^\infty
\eu^{-r^2}r^{n-1}\dd r = \frac12\int_0^\infty
t^{\frac n2 - 1}\eu^{-t}\dd t = \frac12\Gamma(\frac n2)$. ومنه
\[
n\,V_n\int_0^\infty\eu^{-r^2}r^{n-1}\dd r
= V_n\cdot\frac n2\,\Gamma\Bigl(\frac n2\Bigr)
= V_n\,\Gamma\Bigl(\frac n2 + 1\Bigr) = \pi^{n/2} = I_n .
\]
وبما أن الطرفين مبرهنان، يمكن \emph{قراءة} المتطابقة تفكيكا قشريا
لتكامل غاوس: فالكرة ذات نصف القطر $r$ تحمل \omterm{def:b2:multint:domain}{المساحة} $nV_nr^{n-1}$،
ويُكامل الوزن الغاوسي $\eu^{-r^2}$ على القشور.

\textbf{14.} $s_0 = V_1 = 2$ (فالكرة $0$ نقطتان)، $s_1 = 2V_2 = 2\pi$، $s_2 = 3V_3 = 4\pi$،
$s_3 = 4V_4 =
2\pi^2$؛ و$\int_0^R s_{n-1}r^{n-1}\dd r = V_nR^n =
v_n(R)$: أي إن \omterm{def:b2:multint:domain}{المساحة} هي المشتقة القطرية للحجم.

\textbf{15.} من أجل $n = 2k$، تعطي صيغة ستيرلنغ (\cref{thm:b2:comparison:stirling}) المقدار
$k! \sim
\sqrt{2\pi k}\,(k/\eu)^k$، ومنه
\[
V_{2k} = \frac{\pi^k}{k!}
\sim \frac{(\pi\eu/k)^k}{\sqrt{2\pi k}}
= \frac1{\sqrt{\pi n}}\Bigl(\frac{2\pi\eu}n\Bigr)^{n/2}
\qquad (n = 2k).
\]
ومن أجل $n$ الفردي: $V_{2k+1} = 2W_{2k+1}V_{2k} \leq 2V_{2k}$، ومنه تتحقق حصور التلاشي فوق الهندسي
نفسها (إلى غاية عامل $2$ وإزاحة بواحد في الأسّ) --- فمن أجل كل
$q >
0$، $V_n = o(q^n)$.

\textbf{16.} $V_2/4 = \pi/4 \approx 0.785$؛ $V_3/8 = \pi/6
\approx 0.524$؛ $V_{10}/2^{10} = \frac{\pi^5}{120\cdot1024}
\approx 0.0025$. وبوجه عام $\frac{V_n/2^n}{V_{n-2}/2^{n-2}}
= \frac{2\pi}{4n} = \frac{\pi}{2n} \to 0$: فالنسبة
تؤول إلى $0$ (تلاشيا فوق هندسي). ومنه فإن الكرة المصمتة المرسومة
داخله تشغل نسبة متلاشية: فحجم المكعب يهاجر إلى زواياه.

\textbf{17.} بحسب السؤال 1، يكون \omterm{pb:b2:multint:1}{حجم الكرة المصمتة} الداخلية ذات نصف
القطر $1 -
\varepsilon$ هو $V_n(1-\varepsilon)^n$، ومنه تحمل القشرة الخارجية النسبة $1 - (1 - \varepsilon)^n \to 1$.
ومن أجل $n = 100$، $\varepsilon = 0.01$: $(0.99)^{100} =
\eu^{100\ln0.99} \approx \eu^{-1.005} \approx 0.366$: أي إن نحو $63\%$ من الكرة المصمتة
يقع على مسافة لا تتجاوز $1\%$ من سطحها.

\textbf{18.} $(W_n)$ متناقصة، ومنه $W_n \leq W_{n-1} \leq
W_{n-2} = \frac{n}{n-1}W_n$: وبالحصر، $W_{n-1}/W_n \to 1$. وبالضرب
في $W_nW_{n-1} = \frac\pi{2n}$: $W_n^2 \sim
\frac\pi{2n}$، أي $W_n \sim \sqrt{\pi/(2n)}$. والحصر السفلي: $W_n^2 \geq W_nW_{n+1} = \frac{\pi}{2(n+1)}$، ومنه $W_n \geq
\sqrt{\pi/(2(n+1))}$
من أجل كل $n$.

\textbf{19.} البسط: يعطي $1 - x^2 \leq \eu^{-x^2}$ المقدار $(1 -
x^2)^{\frac{n-1}2} \leq \eu^{-(n-1)x^2/2}$، ومع $a =
\frac{n-1}2$،
\[
\int_\delta^1(1 - x^2)^{\frac{n-1}2}\dd x
\leq \int_\delta^\infty\eu^{-ax^2}\dd x
\leq \int_\delta^\infty\frac x\delta\,\eu^{-ax^2}\dd x
= \frac{\eu^{-a\delta^2}}{2a\delta}
= \frac{\eu^{-(n-1)\delta^2/2}}{(n-1)\delta}.
\]
والمقام: $2W_n \geq \sqrt{2\pi/(n+1)}$ بحسب السؤال 18. والقسمة تعطي الحصر المكتوب. ومن أجل
$\delta =
s/\sqrt{n-1}$ يصير $\sqrt{\tfrac{n+1}{2\pi(n-1)}}\;\eu^{-s^2/2}/s = O\bigl(
\eu^{-s^2/2}/s\bigr)$، بانتظام بالنسبة إلى $n$: فخارج الشريحة
$\abs{x_1} \leq s/\sqrt{n-1}$ لا يكاد يوجد حجم، من أجل $s$ كبيرة نسبيا --- وبالتناظر
يصح الأمر نفسه في كل اتجاه.

\textbf{20.} يتعايش النصان لأنهما يصفان إحداثيات مختلفة للنقطة
نفسها. فكل نقطة من $B_n(1)$ تقريبا معيارها قريب من $1$ (السؤال
17: التركيز القطري قرب الكرة)، ومع ذلك تكون كل إحداثية من إحداثياتها
$n$ صغيرة، من رتبة $1/\sqrt n$ (السؤال 19)، وهو أمر منسجم لأن $n$
إحداثية حجمها $1/\sqrt n$ يكون معيارها من رتبة $1$. فالحجم العالي
البعد يتركز حيث تتقاسم جميع الإحداثيات ميزانية \omterm{def:b2:nvs:norm}{المعيار} بالتساوي ---
أي قرب الكرة، لكن بعيدا عن كل قطب محور إحداثي.

\textbf{21.} نشرّح $\Delta_n$ عند $x_n = t \in \intcc01$: فالمقطع هو $\{x' \in \R^{n-1} : x_i \geq 0,\ \sum x_i \leq 1 -
t\} = (1-t)\Delta_{n-1}$، وحجمه
$(1-t)^{n-1}\operatorname{vol}(\Delta_{n-1})$ بالتجانس. ومنه
\[
\operatorname{vol}(\Delta_n) =
\operatorname{vol}(\Delta_{n-1})\int_0^1(1 - t)^{n-1}\dd t =
\frac{\operatorname{vol}(\Delta_{n-1})}{n}
\quad\Longrightarrow\quad
\operatorname{vol}(\Delta_n) = \frac1{n!}\,.
\]

\textbf{22.} تقطع أثمان الإشارات وعددها $2^n$ المجموعة $C_n$ إلى
$2^n$ نسخة من $\Delta_n$ (مع تراكبات مهملة على المستويات الزائدة
الإحداثية): $\operatorname{vol}(C_n) = \frac{2^n}{n!}$. فإذا كان $\sum\abs{x_i} \leq 1$ فإن $\sum x_i^2 \leq
\bigl(\sum\abs{x_i}\bigr)^2 \leq 1$: $C_n \subseteq B_n(1)$؛ و$B_n(1) \subseteq \intcc{-1}1^n$
لأن $\abs{x_i} \leq
\norm x$. ومنه $\frac{2^n}{n!} \leq V_n \leq 2^n$ --- وهو منسجم مع السؤال 15، الذي يضع $V_n$
بين سلّمي العاملي والهندسي.

\textbf{23.} من أجل $(x, y)$ في قرص الوحدة، يكون مقطع $B_4(1)$ هو القرص
ذا نصف القطر $\sqrt{1 - x^2 - y^2}$ في المستوي $(z, w)$، ومساحته $\pi(1 - x^2 - y^2)$. وبالإحداثيات
القطبية:
\[
V_4 = \iint_{x^2+y^2\leq1}\pi(1 - x^2 - y^2)\,\dd x\,\dd y
= \pi\int_0^{2\pi}\!\!\int_0^1(1 - \rho^2)\rho\,
\dd\rho\,\dd\alpha
= \pi\cdot2\pi\cdot\frac14 = \frac{\pi^2}2 ,
\]
وهو يوافق السؤال 6.

\textbf{24.} الاحتمال هو نسبة الحجمين $\dfrac{V_{20}}{2^{20}} = \dfrac{\pi^{10}}{10!\cdot2^{20}}
\approx \dfrac{0.0258}{1\,048\,576} \approx
2.5\cdot10^{-8}$. وعدد السحبات حتى أول
إصابة من رتبة المقلوب، أي نحو $4\cdot10^7$: فمعاين بالرفض كان يعمل بشكل
رائع من أجل القرص (بنسبة إصابات $\pi/4$) يصير عديم الفائدة في البعد
$20$ --- وهي لعنة الأبعاد في سطر واحد.

\textbf{25.} استعمل الطريق الأول (الجزءان الأول والثاني): التشريح من
نوع فوبيني للتكامل المكرر، والتعويض ذا المتغير الواحد في كل إحداثية
(بالتجانس)، \omterm{pb:b2:multint:1}{وتكاملات واليس} --- أي حساب \omterm{def:b2:diffcalc:differential}{تفاضل} وتكامل بمتغير واحد
خالص مضافا إليه التراجع. واستعمل الطريق الثاني (الجزء الثالث):
مبرهنة فوبيني من أجل البنية الجدائية للمقدار $I_n$، وتغيير
المتغيرات القطبي عبر تكامل غاوس في \cref{ex:b2:multint:polar}، والمعادلة الدالية للدالة
$\Gamma$. ويلتقيان في $V_n =
\pi^{n/2}/\Gamma(\frac n2 + 1)$، مع تحويل صيغة ستيرلنغ (\cref{thm:b2:comparison:stirling}) الصيغة
إلى سلوك مقارب. ويعيد مجلد السنة الثالثة بناء هذا كله على تكامل
لوبيغ: فهناك تكون مبرهنة فوبيني وتغيير المتغيرات مبرهنتين من أجل
الدوال القابلة للمكاملة العامة، وتوجد الإحداثيات الكروية في كل بعد،
وتعود حجوم الكرات المصمتة نفسها عوائد محسوبة في مسألتي قياس الجداء
وستيرلنغ --- مع إحلال التقارب المهيمن محل حصورنا اليدوية.
\end{solution}
